\chapter*{Ringraziamenti}
\addcontentsline{toc}{chapter}{Ringraziamenti}
\markboth{Ringraziamenti}{Ringraziamenti}

Desidero innanzitutto esprimere la mia più profonda gratitudine al mio relatore, il \textbf{Prof. Franco Milicchio}, per avermi accordato la fiducia necessaria ad approfondire un argomento per me di straordinaria importanza. La sua guida mi ha permesso di lavorare a un progetto che sento profondamente mio; non avrei potuto desiderare un'esperienza di tesi più gratificante e stimolante. 

\vspace{1em}

\noindent
Un ringraziamento speciale è dedicato alla colonna sonora che ha accompagnato le mie innumerevoli ore di ricerca e di studio. La musica è stata un rifugio, una valvola di sfogo e una fonte inesauribile di energia e concentrazione nei momenti più complessi. Per questo, desidero esprimere la mia profonda gratitudine ai miei artisti preferiti, le cui note sono state una presenza costante e fondamentale in questo percorso: \textbf{Martijn Gerard Garritsen}, \textbf{Michaël Bella}, \textbf{Melle Stomp}, \textbf{Massimo Pezzali}, \textbf{Olivier J. L. Heldens}, \textbf{Julian Dobbenberg}, \textbf{Emilio Justin Behr}, \textbf{Pierre David Guetta}, \textbf{Robbert van de Corput}, \textbf{Timothy Jude Smith}, \textbf{Christopher Comstock}, \textbf{Vladimir Nicolaj de Groot}, \textbf{Aleksandr} e \textbf{Yury Parkhomenko}, \textbf{Jonnie Macaire}, \textbf{Harry Bass}, e \textbf{Victor} e \textbf{Stephan Leicher}. Grazie a tutti voi per avermi fornito la carica necessaria per arrivare fino a questo traguardo.

\vspace{1em}

\noindent
Guardandomi indietro, dai primi anni sui banchi di scuola fino alle aule universitarie, mi rendo conto della fortuna immensa che ho avuto nell'incontrare amicizie che un tempo non avrei nemmeno osato sperare di trovare. Questo traguardo è anche vostro, perché senza le risate, i consigli e i momenti condivisi, questo percorso sarebbe stato infinitamente più difficile. Per questo, voglio ringraziare uno a uno chi ha reso tutto questo più bello:

Un ringraziamento immenso e carico di nostalgia va a tutta \textbf{JEEP PARMA NL}. Se mi fermo a guardare indietro, la mente corre subito a noi, vicini di banco al liceo. Sembra ieri che la nostra più grande e unica preoccupazione era cercare di non scoppiare a ridere semplicemente incrociando i nostri sguardi nel bel mezzo delle lezioni. Ne è passato di tempo da quelle mattine leggere, siamo cresciuti e le aule sono cambiate, ma la nostra complicità è rimasta la stessa di allora. Questo traguardo non è solo mio, è dedicato profondamente anche a voi, i fratelli che quegli anni mi hanno regalato. Grazie al \textbf{Dott. Pietro Begnardi}, a \textbf{Jacopo Mazzoli}, a \textbf{Elia Marchini}, al \textbf{Dott. Niccolò Tuttini} ed al \textbf{Dott. Leonardo Raspanti}. Grazie per essermi stati vicini anche se lontani.

\noindent
Un ringraziamento speciale va a tutta \textbf{EIGHTMATE}. Potrà sembrare insolito ringraziare i propri compagni di giochi in una tesi, ma per me è assolutamente doveroso. Siete stati la mia valvola di sfogo per eccellenza, l'ancora che mi ha permesso di staccare la spina quando il peso degli studi si faceva insopportabile. Grazie a \textbf{Emilio Meini}, \textbf{Michael Manca}, \textbf{John Gimutao}, \textbf{Claudio Gianneschi} e \textbf{Alessandro Trafeli} per le innumerevoli serate passate in cuffia, per le risate che hanno spazzato via lo stress degli esami e per aver reso questi anni decisamente più leggeri. Siete stati una parte fondamentale per questa mia avventura.

Un ringraziamento speciale va al \textbf{Dott. Gabriele Mutatempo}, il mio primissimo amico qui a Roma. Sei stato al mio fianco letteralmente dal giorno zero di questo percorso universitario, diventando subito una certezza in una città nuova. Ammiro profondamente la persona che sei: un amico eccellente su cui poter sempre contare e un fidanzato fenomenale, la tua presenza in questi anni è stata semplicemente fondamentale. Grazie per avermi aperto le porte del tuo mondo e per avermi fatto sentire subito parte del tuo fantastico gruppo di amici, a cui va il mio più sincero affetto: un grazie di cuore, quindi, alla \textbf{Dott.ssa Clara Tondi}, alla \textbf{Dott.ssa Ariane Pepe}, al \textbf{Dott. Tommaso Mandolini}, a \textbf{Tommaso Maniscalco}, alla \textbf{Dott.ssa Francesca Mercante} e al \textbf{Dott. Nicola Tommaso Salerno}.

Non potevo non dedicare un paragrafo al \textbf{Dott. Giacomo Ghinelli}, mio compagno di avventure fin da quella seconda classe superiore passata insieme a Gabriele. Anche se la vita e l'università ti hanno portato a studiare a Milano, la distanza non ha minimamente scalfito la nostra amicizia. Voglio augurarti di conquistare ogni tua meta e di realizzare tutti i tuoi sogni: so perfettamente chi sei e non ho alcun dubbio che ci riuscirai brillantemente. Ti voglio bene, amico mio.

\vspace{1em}

\noindent
Affrontare un percorso universitario in una grande città può spaventare, ma ho avuto la fortuna di incontrare persone che hanno saputo trasformare Roma in una vera e propria casa. A tutti gli amici dell'università, a chi ha condiviso con me infinite pause caffè, nottate passate sugli schermi e ansie prima degli appelli: grazie dal profondo del cuore. Avete reso questi anni un'avventura straordinaria, costruendo con me una famiglia accademica che porterò sempre nel mio cuore.

Grazie a \textbf{Francesco De Marco}. Sei un amico raro, una di quelle persone sempre pronte ad ascoltare nei momenti di crisi e a offrire supporto senza esitazioni. Ti auguro di raggiungere tutti i tuoi traguardi futuri con la stessa ostinata dedizione con cui cerchi la build perfetta a Balatro. Che la vita ti riservi sempre i Joker migliori e moltiplicatori altissimi.

Grazie a \textbf{Diana Buffo}. Grazie per avermi ascoltato nei momenti di crisi, per i consigli preziosi e per la pazienza infinita nell'aiutarmi, anche quando significava ripetermi le stesse cose cento volte. Un grazie speciale va ovviamente anche ad Alca, che palesemente vuole molto più bene a me che a te, e che prima o poi sarà mia.

Grazie a \textbf{Stefano Pedaletti}. Pur essendo l'ultimo arrivato nel gruppo, sei stato subito una presenza fondamentale. Sei il capo della Landa degli Evocatori, anche se le tue urla in chat vocale per le ranked perse dicono esattamente il contrario. Grazie per essere sempre stato al mio fianco, anche quando non volevo, come in quel vicolo di Roma con Nazar. Ricordati sempre, SEI MATTO. Ma ti voglio bene lo stesso.

Grazie a \textbf{Mattia Ambrosio}. Si dice che nessuno dovrebbe incontrare il proprio eroe. Mattia mi ha conosciuto il primo giorno di università e da allora non ci siamo mai separati. Scherzi a parte, sei stato un nucleo essenziale del quartetto iniziale e un amico prezioso con cui condividere risate, sfoghi e pianti. Grazie per essere stato al mio fianco in questo percorso e per aver reso tutto più leggero con la tua presenza... il che compensa in parte la furia che mi infondi quando fai quelle smirate allucinanti su Rainbow Six.

Grazie a \textbf{Daniele Di Lorenzo}. Sei la mia risata quotidiana, una persona a cui è letteralmente impossibile non voler bene. Non dimenticherò mai l'orale di Fondamenti di Telecomunicazioni e quel nostro abbraccio di silenziosa, liberatoria esultanza nei corridoi, consapevoli di avercela finalmente fatta. Certo, per me rimane un mistero assoluto come tu riesca a passare esami del genere essendo fermamente convinto che un proiettile sparato in aria finisca dritto nello spazio, ma sorvoliamo. Grazie per esserci stato in quel momento e in tanti altri, condividendo le sfide, le vittorie e tutte le assurdità di questo percorso.

Grazie a \textbf{Nazar Romanovych Sprynskyy}. Sei senza dubbio una delle persone più brillanti che conosca, al punto che sono quasi certo tu sogni direttamente in Assembly. Grazie per aver sempre condiviso con me le tue infinite conoscenze e per aver eguagliato, fin dal primo giorno, la mia passione per l'informatica, per i videogiochi e per la musica. Sei un amico prezioso e sono grato di averti incontrato, anche se continuo a non capacitarmi di come una mente così geniale possa avere l'algoritmo di Instagram più cursed e inquietante dell'intero emisfero boreale.

\vspace{1em}

\noindent
Un ringraziamento speciale va ai miei compagni "falliti" fuori corso.
Siete stati una presenza costante e un supporto fondamentale per me, siete stati i miei fratelli, per questi 3 anni, e non avrei potuto desiderare compagni di viaggio migliori. Grazie per aver condiviso con me le sfide, le difficoltà e le vittorie di questo percorso, e per aver reso tutto questo più bello e divertente. Siete stati una parte importante di questa avventura, e non vedo l'ora di continuare a condividere con voi i prossimi capitoli della nostra vita accademica.

Grazie a \textbf{Tommaso Tabacchini}. 
Forse dovrei chiamarti con il tuo vero nome TomTab17, o TommyTabby, o TubyTuby, o TibyTiby, o forse, alla fine, semplicemente Tommaso. Se qualcuno ci guardasse da fuori, penserebbe che siamo le due persone più incompatibili sulla faccia della terra. Non ci siamo mai trovati d'accordo su nulla, abbiamo discusso fino a perdere la voce, sbattuto la testa l'uno contro l'altro un'infinità di volte. Eppure, sotto tutte quelle scintille, si nasconde la magia più grande del nostro legame. Ho capito che ci si scontra con questa ferocia solo quando ci si vuole un bene disarmante. Ogni nostra lite non è mai stata un muro, ma un modo per aggrapparci l'uno all'altro, rendendoci inesorabilmente più fratelli. Mi hai insegnato l'onestà brutale di un'amicizia vera, quella che resta anche quando ti volteresti dall'altra parte. Sei stato la mia spina nel fianco preferita e questo viaggio sarebbe stato gelido e vuoto senza di te. Sono immensamente orgoglioso dell'uomo che sei e fiero di averti al mio fianco. Ti voglio bene.

\noindent
P.S. Non illuderti: non vedo l'ora di continuare a darti torto per i prossimi due anni.

Grazie al \textbf{Dott. Emanuele Orlandi}. 
Lele, fatico a trovare le parole giuste, perché come si fa a racchiudere in poche righe la persona che ti ha letteralmente tenuto a galla? La verità, cruda e semplice, è che senza di te, oggi, io non sarei qui a scrivere queste pagine. Sei stato la mia scialuppa di salvataggio nei giorni in cui volevo mollare tutto, quando gli esami sembravano scogli insormontabili, l'ansia mi toglieva il respiro e la fatica minacciava di inghiottirmi. In quei momenti di buio totale, tu sei stato la mia luce. Abbiamo attraversato fianco a fianco l'inferno delle sessioni disperate, i silenzi pesanti dopo una delusione, per poi esplodere di gioia e abbracciarci per ogni singola promozione. Ma oltre ai libri c'è stata la vita vera: la magia dei festival, la musica che ci ha tenuti svegli e vivi, le nottate infinite a confidarci progetti, paure e sogni. Sei stato la mia spalla assoluta, il confidente silenzioso a cui ho affidato le mie fragilità più intime, l'unico che ha sempre saputo come rimettermi in piedi senza mai farmi pesare le mie cadute. Grazie per aver creduto in me quando io per primo non ci riuscivo, e per non avermi lasciato solo nemmeno per un secondo. Indossare questa corona d'alloro ha senso solo perché l'abbiamo conquistata insieme: è tanto mia quanto tua. Ti voglio bene.

\vspace{1em}

\noindent
Il ringraziamento più dolce, è per la mia fidanzata, \textbf{Maddalena Antonacci}, e per la sua famiglia che mi ha sempre fatto sentire a casa. Nana, tu sei la mia fonte di ispirazione quotidiana: mi spingi silenziosamente a fare di più, a puntare più in alto e, soprattutto, a essere un uomo migliore. So che spesso tendi a dubitare di te stessa e a sminuire il tuo valore, ma vorrei che per un attimo potessi guardarti attraverso i miei occhi. Scopriresti una persona meravigliosa di cui sono, e sarò per sempre, infinitamente orgoglioso (anche se devo costantemente competere per il tuo amore con Punch, le donzelle del Peggi, gli smashini, le poké e i canetti). In questi due splendidi anni di amore sei stata la persona che mi è stata più vicina: la tua dolcezza e le tue premure hanno alleggerito i miei momenti di sconforto e hanno moltiplicato la gioia in quelli felici. Grazie per essere rimasta al mio fianco, in ogni singolo istante. Ti amo.

\vspace{1em}

\noindent
Sono senz'altro dovuti altri ringraziamenti, e se mi sono dimenticato di qualcuno, sappiate che non è stato assolutamente intenzionale. Vi voglio bene, e vi ringrazio tutti per aver reso questo percorso così speciale.
Grazie ad \textbf{Alfredo Baldi}, \textbf{Alessandro Fiorella}, \textbf{Ernesto De Salvatore}, \textbf{Lorenzo Conte}, \textbf{Sara Baranes}, \textbf{Mario Savarese}, \textbf{Alessio Paparcurio}, \textbf{Leonardo Centurioni}, il \textbf{Dott. Alessandro Bertoldi}, il \textbf{Dott. Daniel Sadun}, \textbf{Tommaso Cenciotti}, \textbf{Alberto De Paolis}, la \textbf{Dott.ssa Gioia Buratti}, la \textbf{Dott.ssa Lucrezia Dell'Anna}, \textbf{Lorenzo Giannetti}, \textbf{Marina Mozzetta}, il \textbf{Dott. Domenico Avella}, \textbf{Alessandro Bucalossi}, \textbf{David Barbagli}, \textbf{Aldo Giuliani} e la \textbf{Prof.ssa Letizia Maestrucci}.
Chi più e chi meno, chi è qui e chi è lontano, chi è stato al mio fianco per tutto questo tempo e chi è stato una presenza più occasionale, ognuno di voi ha contribuito a rendere questo percorso più ricco, più divertente e più significativo. Grazie di cuore a tutti voi.

\vspace{1em}

\noindent
Un pensiero colmo d'affetto è rivolto a tutta la mia famiglia allargata: ai miei nonni, agli zii e ai cugini. Siete la rete affettiva che ha colorato la mia vita fin da quando ero bambino, una fonte inesauribile di calore, saggezza e incoraggiamento. Nei vostri sguardi sempre orgogliosi ho trovato spesso la spinta per fare un passo in più, con la consapevolezza e la fortuna di avere alle spalle un'intera squadra pronta a fare il tifo per me. Grazie per aver reso ogni tappa di questo percorso un motivo per festeggiare insieme.

Un ringraziamento pieno di dolcezza e venerazione va alle mie nonne, \textbf{Valda Marchi} e \textbf{Annamaria Giannini}. Siete le indiscusse matrone di casa, le radici forti su cui si regge la nostra famiglia. Con la vostra saggezza, i vostri sorrisi e il vostro amore inesauribile, siete state e sarete per sempre le mie guide più preziose. Questo traguardo è anche il frutto dei vostri insegnamenti e del calore infinito che mi avete sempre donato.

Un pensiero speciale, intimo e silenzioso, vola fin lassù ai miei nonni, \textbf{Adriano Conforti} e \textbf{Paolo Turi}. Quanto avrei voluto potervi abbracciare oggi e vedere i vostri occhi brillare di orgoglio per questo traguardo. Anche se le nostre mani non possono stringersi per festeggiare, sento la vostra mancanza in ogni respiro di questa giornata. Spero che da lassù mi stiate guardando e stiate sorridendo. Vi voglio bene e vi penso.

Grazie ai miei zii, \textbf{Guido Conforti} e \textbf{Manuela Conforti}. 
Ci sono legami familiari che diamo per scontati, ma il vostro supporto per me non lo è mai stato. Siete quel tipo di zii che sanno esserci senza mai essere invadenti, che festeggiano i miei traguardi come se fossero i loro e che sanno alleggerire la tensione con un sorriso sincero. Grazie a entrambi.

Grazie ad \textbf{Alessia}. 
Non ci legherà la genetica, ma la famiglia è prima di tutto quella che si sceglie. So che non sei mia zia, ma ti ho sempre considerata come tale, e ti voglio bene come se lo fossi. Grazie per essere sempre stata una presenza dolce e affettuosa nella mia vita, e per essere sempre stata pronta a dispensare consigli preziosi. Grazie davvero.

Grazie alla \textbf{Dott.ssa Rachele De Simone}.
Costruire un legame e far incontrare caratteri diversi non è quasi mai un percorso immediato, ma è proprio nel tempo condiviso che si impara a capirsi e a darsi valore. Ti ringrazio per essere una presenza così centrale e positiva nella vita di mio fratello, e per tutti i momenti vissuti insieme alla nostra famiglia in questi anni. Sono felice di averti qui oggi e ti auguro il meglio per il tuo futuro.

Grazie ai miei cugini: \textbf{Emilio Meini} e \textbf{Sofia Conforti}. 

Emi, all'anagrafe sei mio cugino, ma in cuor mio ti ho sempre considerato a tutti gli effetti come un secondo fratello. Siamo cresciuti insieme ed insieme siamo maturati. Abbiamo condiviso l'infanzia, l'adolescenza e ora stiamo condividendo anche questo. Sono orgoglioso di te, di tutto quello che hai raggiunto e di tutto quello che raggiungerai. Grazie per esserci sempre stato, ti voglio bene.

A Sofia. Continuo a considerarti la mia "cuginetta" per affetto, ma guardandoti oggi mi rendo conto che non lo sei più: sei diventata una giovane donna straordinaria e, ammettiamolo, la mente più brillante e intelligente di tutta la famiglia. Sono impaziente e curioso di vedere tutto quello che realizzerai, e sono orgoglioso di essere tuo cugino. Ti voglio bene.

\vspace{1em}

\noindent
Infine, il ringraziamento più grande, profondo e viscerale è custodito per la mia famiglia più stretta. Siete le mie radici in un mondo che cambia velocemente, il mio riparo dove rientrare a riprendere fiato e la bussola infallibile che mi ha orientato nei giorni in cui l'ansia e la fatica mi facevano smarrire la rotta. So bene che le pagine di questa tesi portano la mia firma, ma l'inchiostro con cui sono scritte è fatto dei vostri sacrifici silenziosi, delle vostre rinunce mai fatte pesare e dell'amore incondizionato che mi ha protetto in ogni istante. Se oggi indosso questa corona d'alloro, è perché voi mi avete costruito la strada mattone dopo mattone, asciugando le mie frustrazioni e festeggiando ogni mio piccolo passo avanti. Questa non è solo la mia laurea, è a tutti gli effetti la nostra vittoria. A mio padre, a mia madre e a mio fratello: i tre pilastri indiscussi della mia vita.


\noindent
A mio fratello, il \textbf{Dott. Leonardo Conforti}.
Definirti semplicemente un fratello maggiore sarebbe riduttivo: sei il mio punto di riferimento assoluto e, prima di ogni altra cosa, il mio migliore amico. Sei la persona con cui ho condiviso letteralmente ogni singolo passo della mia vita e che, a sua volta, ha scelto di condividere tutto con me. So bene di non essere stato il fratello minore perfetto, né tantomeno il coinquilino ideale, eppure non mi hai mai negato la tua ala protettiva, coprendomi le spalle sempre, senza la minima esitazione. Insieme siamo diventati inseparabili, l'uno l'estensione dell'altro: l'altra inarrestabile metà del duo CONFO dietro la console, ma soprattutto complici perfetti nella vita di tutti i giorni. Guardandoti, vedo un esempio straordinario di dedizione, perseveranza e successo: sei la persona che mi rende più fiero al mondo, e il mio desiderio più grande, con questo traguardo, è riuscire a rendere fiero te. Il tuo supporto, i tuoi consigli e la tua presenza costante sono stati la mia vera forza. Grazie per essermi sempre stato vicino. Ti voglio bene Daddo.

\vspace{1em}

\noindent
A mia madre, la \textbf{Dott.ssa Elena Turi}.
Mamma, sei da sempre la mia roccia, il mio porto sicuro e la fonte inesauribile di quell’amore assoluto che mi ha permesso di arrivare fin qui. Se chiudo gli occhi e guardo indietro, vedo le tue braccia pronte a tenermi in collo quando ero piccolo, la tua infinita pazienza seduta accanto a me per aiutarmi con i compiti, e le tue notti in bianco passate a vegliare su di me ogni singola volta che stavo male. Grazie per aver fatto da scudo alle mie paure, per avermi sempre difeso a spada tratta e per avermi cresciuto avvolto in una sensazione di protezione totale. Hai sacrificato silenziosamente una parte di te per darmi lo slancio e il coraggio di inseguire i miei sogni, credendo in me anche quando io stesso faticavo a farlo. Non sei solo la mia mamma: sei la mia migliore amica, la mia confidente più preziosa e la mia più grande fan. Questa corona d'alloro, oggi, la mettiamo in testa a entrambi. Non potrei essere più orgoglioso e grato di essere tuo figlio. Ti voglio bene Mamma.

\vspace{1em}

\noindent
A \textbf{Riccardo Conforti}.
Babbo, sei da sempre il mio modello assoluto di forza, saggezza e integrità. Se oggi sono l'uomo che sono, il merito è in gran parte tuo. Grazie per avermi insegnato, con i fatti e con il sudore prima ancora che con le parole, cosa significhi essere una persona onesta, lavoratrice e profondamente rispettosa. Grazie per avermi sempre sostenuto e per avermi coperto le spalle in ogni mia singola scelta, facendo il tifo per me anche quando quelle scelte si rivelavano sbagliate.
Lo sappiamo entrambi: abbiamo due caratteri forti, ci siamo scontrati e abbiamo litigato innumerevoli volte. Per anni ho vissuto questo nostro essere costantemente in conflitto come un problema, come un ostacolo tra noi. Ma oggi, guardandomi indietro con occhi più maturi, mi rendo conto che è stato proprio il fuoco di quegli scontri a forgiarmi. Mi hai insegnato a tenere testa al mondo, a difendere le mie idee e a non abbassare mai lo sguardo. Dietro ogni nostra incomprensione, dietro ogni porta sbattuta, c'è sempre stato un bene immenso e incrollabile. Ti ho sempre voluto bene e te ne vorrò sempre, incondizionatamente.
Grazie per essere mio padre, il mio eroe e il mio più grande sostenitore. Oggi, con questa tesi tra le mani, voglio farti una promessa. Ti prometto che accorcerò le distanze, che alzerò il telefono per chiamarti più spesso, anche solo per sentire come stai. E ti prometto che, d'ora in poi, sarò più spesso io a fare il primo passo e a dirti per primo quella frase che ci unisce da sempre: Tanto bene a babbo.

\vspace{2mm}
\begin{flushright}
    \textit{Edoardo}
\end{flushright}